\section{Related Work}
\paragraph{Training-free reuse and forecasting.}
TeaCache estimates nonuniform output variation from inexpensive, timestep-modulated inputs and rescales the estimated changes for adaptive caching \citep{liu2025teacache}. Its input sensitivity is important: our distinction is not simply that prior work ignores the current latent. TaylorSeer replaces direct reuse with Taylor-series feature forecasts whose derivatives are approximated from historical feature values \citep{liu2025taylorseer}. \method\ instead proposes an explicit state-and-time model of the guided output. Its temporal fallback is an output-level first-order extrapolator, not a reproduction of TaylorSeer's intermediate-feature implementation. Comparisons must use the actual baseline code and identify any model-specific adaptation.

\paragraph{Guidance- and propagation-aware cache control.}
RA-CFGCache identifies both a branch-guided mismatch and a local-final mismatch under CFG \citep{liu2026racfgcache}. It composes branch-wise error proxies using guidance coefficients and calibrated cross-branch alignment, then applies a timestep-dependent propagation prior before online thresholding. Consequently, propagation-aware caching itself is prior work, not our novelty claim. The proposed distinction is a causal, prompt-specific local operator shared by output prediction and finite-horizon action selection. Whether that operator adds information beyond a timestep prior is unresolved and must be measured; substituting such a prior into \method\ is a necessary ablation, not evidence for online identification.

\paragraph{Numerical sampling.}
DDIM and UniPC reduce the number of denoising evaluations through sampling schemes with different update rules and solver histories \citep{song2021ddim,zhao2023unipc}. \method\ keeps the selected sampling grid and scheduler fixed and changes which steps receive exact network outputs. Its propagation model must therefore use the actual scheduler, including prediction-type conversion and multistep history. A one-step tangent formula is not automatically valid for a predictor--corrector scheduler.

\paragraph{Online identification and predictive control.}
Secant and quasi-Newton methods construct derivative approximations from function evaluations without forming a dense Jacobian \citep{broyden1965class}. We borrow this limited-memory identification principle, but fitting along a denoising trajectory is not equivalent to observing fixed-time state perturbations. Model predictive control repeatedly optimizes a finite-horizon action sequence and executes only its first action \citep{rawlings2017mpc}. \method\ uses that structure for compute allocation; low-rank fitting, receding-horizon planning, and residual checks are not individually claimed as new. The research hypothesis is that coupling them through one useful dynamics model yields a better cache controller.

\paragraph{Video models and evaluation.}
CogVideoX and Wan provide the two video backbones used in the planned comparison \citep{yang2025cogvideox,wan2025wan}. VBench supplies a multidimensional video-evaluation framework \citep{huang2024vbench}. Using its prompt resource for paired LPIPS, SSIM, and PSNR is distinct from computing a full VBench score. We follow RA-CFGCache's shared generation protocol for reproducibility and will report these evaluation scopes separately.
